%=============================================================
%  pages/exp1.tex — Experiment 2
%  
%=============================================================

\experiment{2}{%
    Take a dataset and classify the data based on measurement scales(nominal, ordinal, internval, and ratio). Provide appropriate examples for each category.
}

% ── 1. Dataset Info ───────────────────────────────────────────
\subsection{Objective}
The objective of this experiement is to analyze and classify the dataset into four fundamental levels of measurement scale.


\subsection{Classification of Dataset}
\begin{tabularx}{\linewidth}{ccccX}
    \toprule
    & \textbf{Variable} & \textbf{Data Type} & \textbf{Measurement Scale} & \textbf{Justification} \\
    \midrule
    1 & Name  & Categorial & Nominal & no mathematical relation or order \\
    2 & Gender & Categorial & Nominal & Two categories, no relation \\
    3 & Grade  & Categorial & Ordinal & A > B > C > D > E, intervals are unequal \\
    4 & Marks  & Numerical & Ratio & 0 is nil, 90 marks is twice 45 marks \\
    5 & StudyHours  & Numerical & Ratio & 0 means nil \\
    6 & Attendance  & Numerical & Ratio & 0\% means nil \\

    \bottomrule
\end{tabularx}
